\begin{figure*}[!t]
\centering
\includegraphics[width=\textwidth]{figures/fig-behavior-shift.pdf}
\caption{Behavior shift after \method{}. \method{} sustains longer working sessions and shifts the agent's per-step action mix toward verification on SWE-Bench Pro, and toward execution on Terminal-Bench~2 and \mbox{GAIA-2}.}
\vspace{-1em}
\label{fig:steps}
\end{figure*}

\section{Discussion}

% \looseness=-1 We organize the analysis around four questions, asking how behavior changes after optimization (\S\ref{sec:disc-behavior}), how coreset selection shapes evolution (\S\ref{sec:disc-coreset}), whether updates are consistent (\S\ref{sec:disc-consistency}), and how much retrospective analysis contributes (\S\ref{sec:disc-diagnosis}).

\subsection{How does agent behavior change after optimization?}
\label{sec:disc-behavior}

\looseness=-1 Although \method{} creates new skills and tools for the agent, it is not immediately apparent through which mechanisms these updates enable the agent to perform better on future tasks.
To examine this, we visualize the frequency of tool calls and the cumulative success rate with respect to the number of steps taken by the agent.
Specifically, Figure~\ref{fig:steps} plots the cumulative fraction of held-out tasks resolved within a given number of agent steps, and it shows how the action mix of the agent changes over time.
We observe that the performance improvements across all three datasets primarily originate from higher success rates on tasks requiring long horizons.
In contrast, the gains concentrate on long-horizon tasks rather than those completed in fewer steps, most clearly on SWE-Bench Pro.
Furthermore, the optimization process changes the working patterns of the agent.
Consequently, the optimized agent shifts toward relying more heavily on specific types of actions.
For example, on SWE-Bench Pro, the agent verifies its work much more frequently.
This proactive verification appears to account for a large portion of the performance gains on long-horizon tasks.
On Terminal-Bench~2 and \mbox{GAIA-2}, the agent increases its accuracy by actively applying newly developed tools.
The newly added tools are invoked on nearly every held-out task, and on GAIA-2 the new environment helper runs about twenty times per task.
Terminal-Bench~2 is an instructive exception---there the generated scripts are never executed, and the gain instead comes from the procedural playbook written into the instructions, showing that \method{} adapts the harness surface to each environment rather than always adding executable tools.
Complementary to our action-level statistics, token-level attribution over long reasoning chains offers a finer-grained lens on long-horizon agent behavior \citep{pan2026flashtrace}.

\subsection{How does coreset selection shape capability evolution?}
\label{sec:disc-coreset}

\begin{figure}[t]
\centering
\includegraphics[width=\columnwidth]{figures/fig-coreset-selection.pdf}
\caption{Coreset selection on SWE-Bench Pro. \emph{(a)} Where each selector's picks land on the task embedding, with coverage spreading out, difficulty clustering, and \method{}'s DPP balancing both. \emph{(b)} Held-out pass rate of the harness optimized from each coreset. Difficulty or diversity alone trails even random sampling, and only the DPP's combination reaches the top gain.}
\label{fig:coreset-schematic}
\vspace{-1.5em}
\end{figure}


\begin{table}[t]
\caption{Best-of-$N$ harness proposal vs.\ a single sampled candidate. Held-out pass rate of the $N{=}3$ candidates. \emph{Mean} is the expected score under uniform random selection, and \emph{Chosen} is the candidate \method{} deploys.}
\label{tab:bestofn}
\centering
\footnotesize
\setlength{\tabcolsep}{4pt}
\sbox0{%
\begin{tabular}{lcccc}
\toprule
\textbf{Dataset} & \textbf{Mean} & \textbf{Chosen} & \textbf{Std} & \textbf{Lowest} \\
\midrule
\rowcolor{rowAlt}
SWE-Bench Pro & 0.79 & 0.78 & 0.06 & 0.73 \\
\rowcolor{white}
Terminal-Bench~2 & 0.74 & \textbf{0.76} & 0.03 & 0.71 \\
\rowcolor{rowAlt}
GAIA-2 & 0.34 & \textbf{0.37} & 0.03 & 0.32 \\
\bottomrule
\end{tabular}%
}
\ifdim\wd0>\columnwidth
  \resizebox{\columnwidth}{!}{\usebox0}%
\else
  \usebox0%
\fi
\end{table}

\begin{table}[t]
\caption{Ablation of the diagnosis step across all three benchmarks. Each row re-renders full \method{}'s diagnoses with one cue blanked, holding the proposal agent's prompt fixed (\S\ref{sec:disc-diagnosis}). \emph{Raw trajectory} skips diagnosis and shows the proposal agent the raw trajectories directly. Cells report held-out pass rate.}
\label{tab:diagnosis}
\centering
\footnotesize
\setlength{\tabcolsep}{4pt}
\sbox0{%
\begin{tabular}{lccc}
\toprule
\textbf{Variant} & \textbf{SWE Pro} & \textbf{TB~2} & \textbf{GAIA-2} \\
\midrule
\rowcolor{rowAlt}
Full diagnosis & \textbf{0.78} & \textbf{0.76} & \textbf{0.37} \\
\rowcolor{white}
\;\; $-$ self-consistency & 0.56 & 0.75 & 0.27 \\
\rowcolor{rowAlt}
\;\; $-$ self-validation & 0.70 & 0.73 & 0.30 \\
\rowcolor{white}
Raw trajectory & 0.60 & 0.75 & 0.29 \\
\bottomrule
\end{tabular}%
}
\ifdim\wd0>\columnwidth
  \resizebox{\columnwidth}{!}{\usebox0}%
\else
  \usebox0%
\fi
\vspace{-1.5em}
\end{table}

\looseness=-1 We investigate how coreset selection influences the optimization process.
To this end, we compare our DPP-based selection against several ablation strategies.
These variants include selecting tasks solely by difficulty ($\theta=1$), selecting tasks purely to maximize coverage ($\theta=0$), and sampling trajectories randomly.
In addition, we measure the final performance of the optimized harness under each selection strategy.
As the t-SNE projection of task embeddings in Figure~\ref{fig:coreset-schematic} reveals, selecting tasks exclusively by difficulty causes the chosen samples to cluster in a narrow region of the task distribution.
This clustering occurs because the language model judges certain types of tasks as inherently more difficult, and consequently it fails to include other task types in the coreset.
As a result, this strategy yields no meaningful performance improvement after optimization.
Similarly, optimizing solely for coverage also produces suboptimal results.
In contrast, random sampling can occasionally select a trajectory that proves useful for optimization.
These findings suggest that a coreset selection strategy balancing both difficulty and diversity is vital for providing the proper signals to guide harness optimization.

\subsection{Does \method{} produce consistent harness updates?}
\label{sec:disc-consistency}

Because \method{} operates without ground-truth labels for reference, we evaluate whether its optimization outputs are consistent across runs.
This consistency is closely tied to the overall effectiveness of our best-of-$N$ harness proposal.
We examine whether the selection strategy reliably identifies the candidate harness that performs best on downstream tasks.
In this experiment, we measure the test scores of all three generated candidate harnesses rather than just the most preferred one.
Here Table~\ref{tab:bestofn} shows that the generated harnesses exhibit only moderate variance.
Notably, even the lowest-scoring candidate still meaningfully improves agent performance over the baseline.
Furthermore, the best-of-$N$ selection prevents the deployment of poorly performing harnesses.
Specifically, the chosen harness scores higher than the worst candidate across all three benchmarks.
At the same time, the most preferred harness does not invariably coincide with the highest-scoring candidate on the test set, though the selection consistently avoids the worst candidate.
Comparing each candidate's preference score $S_j$ with its held-out pass rate further shows that self-preference is a conservative selector rather than a perfect ranker: on SWE-Bench Pro a less-preferred candidate in fact reached 0.85, above the chosen one's 0.78.
What the mechanism reliably guarantees is safety rather than perfect ranking.
Because an update is deployed only when the agent prefers it over the current harness ($S_j > 0$), the deployed harness never fell below the vanilla baseline on any benchmark; every candidate beats vanilla on SWE-Bench Pro and GAIA-2, and on Terminal-Bench~2 the worst candidate matches vanilla (0.71) while the chosen one reaches 0.76.



\subsection{How much does retrospective analysis contribute?}
\label{sec:disc-diagnosis}

We analyze the contribution of the two signals extracted during retrospective analysis, namely self-validation and self-consistency.
We compare these explicit signals against a more direct approach that provides raw trajectories during optimization and skips the explicit retrospective analysis step.
In addition, we study the individual contribution of each diagnostic signal to the final performance.

To investigate the necessity of the group rollout and diagnostic stages, we conduct an ablation study.
Specifically, we remove the self-validation and self-consistency signals independently, and we subsequently rerun the optimization and evaluation procedures.
We also introduce a raw trajectory baseline.
This baseline bypasses the separate ranking analysis. Instead, it provides the original trajectories directly to the optimization step, asking the agent to analyze the trajectory and propose improvements in a single pass.
Table~\ref{tab:diagnosis} shows that removing either the self-consistency or the self-validation signal consistently degrades final performance across benchmarks.
This result indicates that both signals are highly important for optimizing the harness.
Additionally, full diagnosis outperforms the simplified raw-trajectory baseline on all three benchmarks.
The raw-trajectory baseline grants the proposal agent the same editable surface as \method{}---tools, skills, and instructions---so the gap to full diagnosis (0.60 vs.\ 0.78 on SWE-Bench Pro) reflects the value of the structured retrospective analysis rather than a broader editing surface.
This result suggests that although single-pass trajectory analysis is feasible, the explicit self-validation and self-consistency signals are essential rather than incidental, yielding more reliable improvements across datasets.
Finally, Table~\ref{tab:diagnosis} is a controlled ablation that removes one diagnosis cue while holding the rest of the pipeline on a fixed path ($n{=}1$), so it measures the marginal contribution of each cue; it should not be read as the end-to-end effect of disabling a signal.
To measure run-to-run robustness, we additionally reran the full pipeline end to end twice per ablated strategy, with independently generated candidate sets: every rerun configuration stayed within $\pm$0.01--0.03 of its mean ($0.79 \pm 0.01$ without self-consistency and $0.815 \pm 0.025$ without self-validation on SWE-Bench Pro, $n{=}2$), and in no rerun did performance drop to the controlled-ablation level.
With $n{=}2$ this is indicative rather than conclusive, but it shows that the best-of-$N$ proposal and the remaining signal compensate when a cue is missing, so removing one component does not destabilize the pipeline.

